\clearpage
\fancyhead[L]{\small NOVA ACTA UPSALIENSIS.}
\hypertarget{native-609-lower}{}
\begin{center}\small Göttingische gelehrte Anzeigen. Stück 70. Seite 692..694. 1818 Mai 2.\\[3pt]\rule{.72\linewidth}{.3pt}\end{center}

\emph{Nova Acta regiae societatis scientiarum Upsaliensis.} Vol. VII. Upsala. 394 Seiten in Quart.

Den Anfang dieser Sammlung macht eine von dem verstorbenen Prof. \textsc{Erich Prosperin} nachgelassene Abhandlung über die Bahnen, welche die Nebenplaneten um die Sonne beschreiben würden,
wenn ihre Hauptplaneten plötzlich vernichtet würden, oder auf jene zu wirken aufhörten. Den Umständen nach können diese Bahnen Ellipsen, Parabeln oder Hyperbeln werden, rechtläufig oder rückläufig. Die Gattung des Kegelschnitts wird theils durch die Entfernung des Nebenplaneten von seinem
Hauptplaneten, theils durch den Platz bedingt, wo eben jener bei der vorausgesetzten Vernichtung des
Hauptplaneten sich in seiner Bahn befindet. Wenn man, mit dem Verf., die Ellipticität der Bahn des
Hauptplaneten um die Sonne und die der ursprünglichen Bahn des Nebenplaneten um seinen Hauptplaneten vernachlässigt, und das Product der Entfernung des Hauptplaneten von der Sonne in die Masse
von jenem (die der Sonne zur Einheit genommen) $= d$ setzt, die Entfernung des Nebenplaneten vom
Hauptplaneten hingegen $= D$, so wird die neue Bahn des Nebenplaneten um die Sonne, nach der Vernichtung des Hauptplaneten, immer eine Ellipse sein, wenn $D$ grösser ist als $(3+\BareRoot{8})d$; sie wird immer eine Hyperbel werden, wenn $D$ kleiner ist als $(3-\BareRoot{8})d$; zwischen beiden Grenzen kann sie Ellipse, Parabel oder Hyperbel werden. Ist $D$ grösser als $d$, so kann die neue Bahn nicht rückläufig
werden; ist $D$ grösser als $d\BareRoot{\tfrac{1}{2}}$, so kann sie wenigstens keine rückläufige Hyperbel werden; endlich ist
$D$ kleiner als $d\BareRoot{\tfrac{1}{8}}$, 
so kann sie keine rechtläufige Ellipse bleiben. Bei diesen Sätzen ist vorausgesetzt,
dass die Masse des Hauptplaneten gegen die der Sonne sehr klein ist, und dass die ursprüngliche Bewegung des Nebenplaneten in der Ebene des Hauptplaneten geschieht; ohne diese Voraussetzungen bedürfen jene einiger Modificationen. Hr. \textsc{Prosperin} hat alle einzelnen Nebenplaneten unsers Sonnensystems besonders betrachtet und sich die Mühe gegeben, sehr ausgedehnte Tabellen dafür zu berechnen.
Die Bahn unsers Mondes würde, wo auch immer er die Erde plötzlich verlöre, eine rechtläufige Ellipse
bleiben; die Jupiters-, Saturns- und Uranustrabanten hingegen könnten den Umständen nach auch parabolische oder hyperbolische Bahnen erhalten. Offenbar lassen sich diese Sätze nun auch umkehren; \SourceLeaf{610}ein in einer parabolischen oder hyperbolischen Bahn in die Nähe des Jupiter, Saturn oder Uranus gelangender Weltkörper, würde um den Planeten eine kreisförmige Bahn antreten, wenn die Einwirkung des Planeten erst nach schon geschehener Annäherung \emph{mit einem Male} anfinge. Da dies aber nicht der Fall der Natur ist, so bleibt die ganze Untersuchung eigentlich nichts, als ein mathematisches Spielwerk. Jenen Umstand hat der Verf. zwar nicht ganz übersehen, aber nicht genug beherzigt, denn es ist gerade eben so \emph{unmöglich}, dass durch die blossen Anziehungskräfte ein zufällig in die Nähe eines Hauptplaneten kommender Comet zu einem Trabanten desselben werden könne, als dass ein schon vorhandener Nebenplanet seinen Hauptplaneten verlasse, um in einer parabolischen oder hyperbolischen Bahn davon zu gehen.
